\section{The model}

This section gives the LaTeX working version of the self-employment model.
It is intended to be the collaboration-safe drafting surface for model writing,
separate from the actively edited LyX file.

\subsection{Environment and timing}

Consider a Lucas (1978) span-of-control environment in which households differ
in human capital $h$. As in Gomez and Kuhn (2017), households are grouped into
three education categories, primary $(P)$, secondary $(S)$, and tertiary $(T)$.
Higher education shifts the distribution of human capital upward and is
positively correlated with managerial ability $x$, worker productivity $z$, and
job-separation risk.

At the start of each period, a household enters as unemployed $(i_w = 0)$,
employed $(i_w = 1)$, or entrepreneur $(i_w = 2)$. The household then chooses
whether to operate a business, remain in work, or search for employment.
Employment is state-dependent: a currently employed worker faces separation
probability $\rho(h)$, while a non-employed worker can only reach employment
through matching, with probability $p(\theta)$. Because current employment
provides an incumbent match, the value of employment is higher for an employed
worker than for an otherwise identical non-employed worker.

The model distinguishes two entrepreneurial margins. First, households differ
between opportunity entrepreneurship and necessity entrepreneurship. Opportunity
entrepreneurs would choose entrepreneurship even if they began the period with a
viable employment relationship. Necessity entrepreneurs choose entrepreneurship
only when they lack that employment option. Second, within entrepreneurship the
business can be either nonemployer self-employment or employer
entrepreneurship. The self-employment branch is designed to quantify both
composition margins jointly.

\paragraph{Timing.}

Occupational choice occurs at the start of period $t$. Households that choose
entrepreneurship operate a business during the period and choose capital and
hired labor. If hired labor satisfies $n_h = 0$, the business is nonemployer
self-employment. If $n_h > 0$, the business is employer entrepreneurship.
Entrepreneurial businesses face closure risk $\rho_e(h)$. In the benchmark,
closure preserves current-period operating income, but moves the household into
non-employment for the next period without automatic regular unemployment
insurance eligibility.

Previously employed workers may continue in work but face separation
probability $\rho(h)$. If separation occurs, they receive unemployment support
for the remainder of the period and do not search again until the next period.
Previously unemployed workers search for employment and match with firms with
probability $p(\theta)$. Successful matches produce labor income within the
period; unsuccessful search leads to non-employment and benefit receipt.

Unlike a standard ex post wage-bargaining search model, the wage here is taken
as given when households make occupational choices. The occupational decision
therefore determines whether a household enters entrepreneurship or directed job
search before the employment realization is known.

\subsection{Preferences, endowments, and technology}

\paragraph{Preferences.}

Households derive utility from consumption $c_t$ with period utility
$U(c_t)$ and discount factor $\beta$.

\paragraph{Endowments.}

Each household begins the period with assets $a$ and an education-linked human
capital draw $h$. Human capital is positively correlated with managerial talent
$x$, worker productivity $z$, and separation risk.

\paragraph{Production.}

Entrepreneurs operate a single-project technology using capital $k$, owner
labor $\ell_e$, and hired labor $n_h$:
\begin{equation}
    y = x k^{\alpha} \left( \ell_e + \xi z n_h \right)^{\gamma},
    \qquad \alpha, \gamma > 0.
    \label{eq:model_prod_self_employment}
\end{equation}
Self-employment is feasible at the corner $n_h = 0$, while employer
entrepreneurship corresponds to $n_h > 0$.

\paragraph{Factor prices and policy.}

Capital rents at market rate $r$, labor is paid wage $w$, and the government
provides unemployment insurance financed through taxation. Let $b$ denote the
unemployment benefit and $\tau$ the labor-tax object borne jointly by workers
and employers, with employer share $\psi$.

\subsection{Static firm problem}

In the self-employment version, the entrepreneur chooses capital and hired labor
while supplying owner labor internally. Let $f_{hire}$ denote the fixed
employer-threshold cost paid whenever the entrepreneur hires strictly positive
labor. The static problem is
\begin{equation}
    \max_{k,\, n_h \geq 0}
    \; x k^{\alpha} \left( \ell_e + \xi z n_h \right)^{\gamma}
    - wz(1 + \psi \tau) n_h
    - rk
    - \kappa \frac{\rho z n_h}{q(\theta)}
    - f_{hire}\mathbf{1}\{n_h > 0\}.
    \label{eq:static_firm_problem_self_employment}
\end{equation}

The term involving $\kappa$ captures the smooth hiring or replacement cost from
search frictions, while $f_{hire}$ captures a discrete hurdle at the first
worker margin. This is the key object that separates nonemployer self-employment
from employer entrepreneurship in the benchmark extension.

For an interior employer choice, the first-order condition for hired labor is
\begin{equation}
    \gamma x k^{\alpha}
    \left( \ell_e + \xi z n_h \right)^{\gamma - 1} \xi z
    = w(1 + \psi \tau) + \frac{\rho \kappa}{q(\theta)}.
    \label{eq:foc_hired_labor_self_employment}
\end{equation}

This implies the policy rule
\begin{equation}
    n_h^*(k,x,\widetilde{w})
    =
    \max\left\{
        0,
        \frac{1}{\xi z}
        \left[
            \left(
                \frac{\gamma x k^{\alpha}\xi z}
                {w(1+\psi\tau)+\frac{\rho\kappa}{q(\theta)}}
            \right)^{\frac{1}{1-\gamma}}
            - \ell_e
        \right]
    \right\}.
    \label{eq:nh_policy_self_employment}
\end{equation}

Substituting the labor policy into the static objective gives entrepreneur
profits
\begin{equation}
    \pi(k,n_h^*,x;\widetilde{w})
    =
    x k^{\alpha}\left(\ell_e + \xi z n_h^*\right)^{\gamma}
    - z n_h^* w(1+\psi\tau)
    - rk
    - \kappa \frac{\rho z n_h^*}{q(\theta)}
    - f_{hire}\mathbf{1}\{n_h^*>0\}.
    \label{eq:profit_self_employment}
\end{equation}

Under perfect capital markets, the interior capital choice satisfies
\begin{equation}
    \alpha x k^{\alpha - 1}\left(\ell_e + \xi z n_h\right)^{\gamma} = r.
    \label{eq:foc_capital_self_employment}
\end{equation}

\subsection{Household problems}

\paragraph{Entrepreneurs.}

Entrepreneurs maximize expected discounted utility
\begin{equation}
    \max_{\{c_t, a_{t+1}, k_t, n_{h,t}\}}
    \mathbb{E}\sum_{t=0}^{\infty}\beta^t U(c_t)
\end{equation}
subject to
\begin{align}
    c + a'
    &\leq
    (1 + r - \delta)a
    + \pi(k,n_h,x;\widetilde{w})
    - \tau_{\pi},
    \label{eq:budget_entrepreneur_self_employment}
\end{align}
and the borrowing constraint
\begin{equation}
    a' \geq -\bar{a}.
    \label{eq:borrowing_constraint_self_employment}
\end{equation}

\paragraph{Employed workers.}

Employed workers solve
\begin{equation}
    \max_{\{c_t, a_{t+1}\}}
    \mathbb{E}\sum_{t=0}^{\infty}\beta^t U(c_t)
\end{equation}
subject to
\begin{align}
    c + a'
    &\leq
    (1+r-\delta)a
    + (1-\rho)\left[wz(1-(1-\psi)\tau_y)\right]
    + \rho b,
    \label{eq:budget_employed_self_employment}
\end{align}
with $a' \geq -\bar{a}$.

\paragraph{Unemployed workers.}

Unemployed workers solve
\begin{equation}
    \max_{\{c_t, a_{t+1}\}}
    \mathbb{E}\sum_{t=0}^{\infty}\beta^t U(c_t)
\end{equation}
subject to
\begin{align}
    c + a'
    &\leq
    (1+r-\delta)a
    + p(\theta)\left[wz(1-(1-\psi)\tau_y)\right]
    + (1-p(\theta))b,
    \label{eq:budget_unemployed_self_employment}
\end{align}
with $a' \geq -\bar{a}$.

\subsection{Government}

The government finances unemployment insurance and any consumption-floor
transfers through taxation. In the endogenous-tax experiments, the relevant tax
object adjusts to satisfy the model's budget-closure rule. In the fixed-tax
experiments, the tax rate is held fixed and transfer changes show up directly in
the private-side outcomes.

\subsection{Occupational choice and equilibrium}

Given prices and policy, the household compares the value of entrepreneurship,
employment, and non-employment. Necessity entrepreneurship is defined as the
set of households that choose entrepreneurship only when they lack an active
employment relationship. Opportunity entrepreneurship is defined as the set of
households that choose entrepreneurship regardless of whether they begin the
period employed or non-employed.

An equilibrium consists of
\begin{itemize}
\item household policy rules for consumption, saving, occupation, capital, and
hired labor,
\item value functions for entrepreneurial, employed, and unemployed states,
\item prices $(r,w,\theta)$ consistent with capital-market and labor-market
clearing,
\item firm policies consistent with the static problem and search costs,
\item and government policy objects consistent with the budget rule used in the
experiment.
\end{itemize}

In the self-employment benchmark, the central equilibrium object is not just the
overall entrepreneur share, but the composition of entrepreneurs between
nonemployer self-employment and employer entrepreneurship, overall and by
education group.
